%!TEX program = lualatex
% ==========================================================================
%  The Natarajan Dichotomy as One D-Sector Readout in Two Regimes:
%  A τ-Categorical Pre-Registration against the Natarajan, Chiang & Dutra
%  (2026, ApJL 1001, L12) Multiscale Cluster-Lensing CDM Crisis.
%
%  Thorsten Fuchs & Anna-Sophie Fuchs
%  Panta Rhei Research Program
%  May 2026
% ==========================================================================
\documentclass{pantarhei-researchnote}

\input{tau-macros}

% ---- PDF metadata (arXiv preflight: merges with class-level \hypersetup) --
\hypersetup{
  pdftitle  = {The Natarajan Dichotomy as One D-Sector Readout in Two Regimes:
               A tau-Categorical Pre-Registration against the Natarajan, Chiang
               and Dutra (2026) Multiscale Cluster-Lensing CDM Crisis},
  pdfauthor = {Thorsten Fuchs and Anna-Sophie Fuchs},
  pdfsubject= {astro-ph.CO; cluster lensing; capacity-gradient gravity;
               Category tau; doi:10.17605/OSF.IO/BFRX3},
  pdfkeywords={dark matter; gravitational lensing; galaxy clusters;
               cluster substructure; capacity-gradient gravity; Category tau;
               categorical pre-registration; Natarajan dichotomy;
               doi:10.17605/OSF.IO/BFRX3}
}

% ---- Note-local macros ----------------------------------------------------
\newcommand{\Nat}{N26}                          % short anchor for Natarajan 2026
\newcommand{\aMOND}{a_0}                        % MOND-effective acceleration scale
\newcommand{\ellcl}{\ell_{\mathrm{cl}}}         % classical-validity scale (host)
\newcommand{\ellsub}{\ell_{\mathrm{cl},\mathrm{sub}}}  % classical-validity scale (subhalo)
\newcommand{\elltauv}{\ell_\tau}                % cosmological tau-length
\newcommand{\Meff}{M_{\mathrm{eff}}}
\newcommand{\Mpa}{M_{\partial}}                 % boundary-holonomy mass
\newcommand{\Mp}{M_{\mathrm{p}}}                % baryonic / particle mass
\newcommand{\Msub}{M_{\mathrm{sub}}}
\newcommand{\Mtwohun}{M_{200}}
\newcommand{\Rtwohun}{R_{200}}
\newcommand{\Mpc}{\mathrm{Mpc}}
\newcommand{\kpc}{\mathrm{kpc}}
\newcommand{\Msun}{M_{\odot}}
\newcommand{\Rs}{R_S}                           % Schwarzschild radius
\newcommand{\rtran}{r_{\mathrm{trans}}}
\newcommand{\rtid}{r_t}
\newcommand{\rcore}{r_{\mathrm{core}}}
\newcommand{\Pggsl}{P_{\mathrm{GGSL}}}
\newcommand{\readout}{\textsc{[readout]}}
\newcommand{\evstag}{\textsc{[evs]}}
\newcommand{\bridge}{\textsc{[bridge]}}
\newcommand{\falsifierID}[1]{\textsc{N16.#1}}

\renewcommand{\arraystretch}{1.2}

\begin{document}

\lane[pr-lane-publications]{Categorical Macrocosm \textperiodcentered\ Book V}
\title{The Natarajan Dichotomy as One D-Sector Readout in Two Regimes}
\subtitle{A $\tau$-Categorical Pre-Registration against the Natarajan, Chiang \& Dutra (2026) Multiscale Cluster-Lensing CDM Crisis}
\noteedition{May 2026}

\author{%
  Thorsten~Fuchs\thanks{\texttt{thorsten@panta-rhei.site}. Corresponding author.}
  \and
  Anna-Sophie~Fuchs%
}

\date{May 2026}

\footerseries{Research Note}
\footerslug{natarajan-cluster-dichotomy-categorical-v1}
\footerdate{May 2026}

% -----------------------------------------------------------------------
\begin{abstract}
\noindent
Natarajan, Chiang \& Dutra~\cite{Natarajan2026} have reported a new flavour
of cold-dark-matter (CDM) crisis from a multiscale strong+weak gravitational
lensing analysis of three massive clusters characterised as dynamically
relaxed in their analysis (MACS~J0416, J1206, J1149; $z = 0.39$--$0.54$;
$\Mtwohun \sim 10^{15}\,\Msun$; J0416's dynamical state is contested in the
broader Hubble Frontier Fields literature), compared to mass-matched
TNG-Cluster ($\Lambda$CDM) analogs. Four substructure diagnostics yield a
\emph{dichotomy}: the same member-galaxy subhalos must carry steep inner
density slopes $\gamma \gtrsim 2.5$--$3$ at $r \lesssim 0.01\,\Rtwohun$
(core-collapsed SIDM-like) \emph{and} collisionless-CDM-like outer tidal
extents at $r \sim 10$--$100\,\kpc$. No single particle-DM microphysics
satisfies all four.
We pre-register a Category-$\tau$ reading of this dichotomy as
\emph{one $D$-sector readout in two regimes}, not two DM components. Three
$\tau$-distinctive positive readings are committed: (i) the outer tidal
extents follow from the Newtonian-limit theorem (V.T28/V.T78) on
tidally-pruned capacity wells at radii where the cluster screening
factor $1 + r/\elltauv \approx 1$ (consistency check passed by construction);
(ii) the projected inner number excess at $R/\Rtwohun \lesssim 0.2$
($5\sigma$--$40\sigma$ in~\cite{Natarajan2026}) is sign-correctly produced
by inner-promotion of the boundary-holonomy mass $\Mpa(R)$ in the host
capacity well, a structurally absent quantity in TNG-C; (iii) the
transition between regimes occurs at
$\rtran \approx \ellsub(\Msub) = \sqrt{G \Msub / \aMOND}$ with
$\aMOND = c H_0 \imark / 2$, predicting $\rtran \approx 3.4$--$10.8\,\kpc$
for $\Msub \in [10^{10}, 10^{11}]\,\Msun$, in order-of-unity agreement
with the empirical break at $r \sim 10\,\kpc$ across all three MACS clusters.
We pre-register one headline falsifier
(\falsifierID{3}: host-mass independence of $\rtran$ at fixed subhalo
mass function, testable with Euclid + Rubin LSST by $\sim 2029$) and one
structurally honest tension reframed as falsifier: the $\tau$-saturation
family predicts \emph{flatter}, not steeper, inner density slopes
(V.P16~\readout); the GGSL probability must therefore be matched by
member-density (this work, R1), not density-slope, with the discriminating
observation deferred to the future GGSL+kinematics inner-mass-content
test. Three derivational gaps are flagged honestly for future Book V
scaffolding: the precise $\gamma$ value, the SHMF amplitude, and the
$5\sigma$--$40\sigma$ inner-excess amplitude.
\end{abstract}

\keywords{dark matter; gravitational lensing; galaxy clusters; cluster
substructure; capacity-gradient gravity; Category $\tau$; categorical
pre-registration; Natarajan dichotomy}

\maketitle

\begin{center}\small\noindent
\textbf{OSF DOI:}\enspace
\href{https://doi.org/10.17605/OSF.IO/BFRX3}{10.17605/OSF.IO/BFRX3}
\enspace\textbar\enspace
arXiv\,(astro-ph.CO\,+\,astro-ph.GA) mirror pending category endorsement
\end{center}
\vspace{0.5em}


% -----------------------------------------------------------------------
\section{Why this paper}\label{sec:why}

The Natarajan~et~al.\ (2026) CDM-crisis paper~\cite{Natarajan2026} differs
from prior small-scale tensions in two respects that matter for any
non-particle response.

\paragraph{Observational robustness.} The $5\sigma$--$40\sigma$ inner
projected substructure excess at $R/\Rtwohun \lesssim 0.2$ survives the
authors' $10^6$-iteration bootstrap injection of force-resolution-disrupted
inner subhalos (their Appendix~B), and the order-of-magnitude tension in
GGSL probability $\Pggsl$ persists under varied baryonic-feedback
prescriptions (their~\S{}9.3, citing
\cite{Meneghetti2020,Meneghetti2023,Ragagnin2022,Heinze2024}).
This is not a marginal anomaly within a noisy substructure inventory; it is
a stress test of $\Lambda$CDM at the dense-cluster intermediate-redshift
regime, performed under maximal $\Lambda$CDM-favourable conditions
(TNG-Cluster, baryonically self-consistent, force-disruption-aware).

\paragraph{Structural distinctness.} Unlike the missing-satellites or
cusp--core problems, the Natarajan tension demands that the \emph{same}
subhalo host two mutually exclusive microphysical signatures:
\emph{collisional-like} inner density slopes
($\gamma \gtrsim 2.5$--$3$ at $r \lesssim 0.01\,\Rtwohun$,
core-collapsed SIDM regime) and \emph{collisionless-like} outer tidal
extents ($\rtid$ at $r \sim 10$--$100\,\kpc$, consistent with NFW dynamical
friction). No single self-interaction cross-section
$\sigma/m$ satisfies both~\cite{Natarajan2026}~\S{}10.

\paragraph{Open door to non-particle alternatives.}
The authors~\cite{Natarajan2026}~\S{}10 explicitly foreclose the
conventional SIDM resolution --- a sophisticated structural commitment
we gratefully acknowledge --- in concert with the
Chiang+2026b~\cite{ChiangSimResolution2026} companion analysis, and
nominate either a hybrid CDM+SIDM model or \emph{``entirely new classes
of DM theories''} as candidate resolutions. This opens the response
landscape to non-particle frameworks --- where the present Panta Rhei
Programme is already positioned via the Category-$\tau$ no-particle-DM
commitment of~\cite{KSZCategorical} (cluster-scale lensing without
particle dark matter, V.T213 universality at $\Meff/\Mp \le 6.65$).
A Category-$\tau$ note can therefore engage substantively, with cluster-scale
precedent already in the public corpus, and without competing inside the
saturating SIDM-core-collapse debate.

What follows is a categorical pre-registration in the
KSZ~\cite{KSZCategorical} and LRD-v3.3~\cite{LRDCategoricalv33} style:
three $\tau$-distinctive positive readings, one numerical anchor, one
headline falsifier, three pre-registered gaps, and one structurally
honest tension reframed as falsifier.

\section{The \texorpdfstring{$\tau$}{tau}-position in one line}\label{sec:tau-position}

\paragraph{No particle DM.} The Panta Rhei Programme commits to a
capacity-gradient framework in which dark-matter phenomenology is a
projection artefact of the categorical $D$-sector readout, not a particle
species. The framing was pre-registered against the Gallardo~et~al.\
pairwise-kSZ force-law constraint~\cite{KSZCategorical}, where the
Newtonian-limit theorem V.T78 fixes the large-scale exponent at
$n = 2$ (Newtonian), in tension with MOND $n = 1$.

\paragraph{KSZ cluster-scale precedent (with substructure caveat).}
\cite{KSZCategorical} commits the framework to cluster lensing
phenomenology via the boundary-holonomy mass closure
$\Meff/\Mp \le 6.65$ (V.D272 + V.T212~\readout), audited against a
five-cluster universality catalogue (V.T213). That catalogue is
\emph{four merging plus one relaxed} (Bullet, A2744, El Gordo,
MACS~J0025; plus A1689), at total-mass per cluster, with $\pm 20$--$50\%$
precision on integrated cluster mass only (not substructure-resolved).
The present three Natarajan clusters are characterised as relaxed
in~\cite{Natarajan2026}, with MACS~J0416's dynamical state contested
in the broader HFF literature, and the observable is substructure-resolved
at $\kpc$ scale. The KSZ V.T213
universality is therefore inherited here only as a calibration prerequisite
on integrated cluster mass, not as a substructure-resolved theorem.

\paragraph{LRD-v3.3 format precedent.}
The pre-registration scaffold follows~\cite{LRDCategoricalv33}, which
attaches eight observational signatures to the heavy-seed black-hole
birth condition at $z \approx 8$--$15$. Signature~3 of that note
($\Meff/\Mp$ aperture-dependent saturation at LRD scales) shares its
anchor (V.D272 + V.T213) with the present~\S\ref{sec:two-regimes} below;
to avoid silent double-booking, the LRD coupling is disclosed in
\falsifierID{3} of the present note (see~\S\ref{sec:falsifiers}).

\paragraph{Status discipline.}
Every $\tau$-statement on Natarajan observables in this note carries
\readout{} (consistency reading or interpretation, not external-physics
theorem) or \evstag{} (external-validation seam, falsifier).
No \textsc{[formal]} or \textsc{[hinge]} claim is made on the
$\gamma$ value, the SHMF amplitude, the $r_t$ value, or the
$5\sigma$--$40\sigma$ amplitude. The V.T203 conjectural status
(cluster capacity mass discrepancy, $D_\tau \sim 2$--$4$ vs.\
$D_{\mathrm{obs}} \sim 5$--$7$~\cite{KSZCategorical} L811--814) is carried
forward verbatim; this note does not narrow that gap, only structurally
locates it (see~\S\ref{sec:gaps}).


% -----------------------------------------------------------------------
\section{Two regimes, two mechanisms}\label{sec:two-regimes}

The Natarajan dichotomy presents at face value as a microphysical paradox:
the same dark-matter sector must be \emph{collisional} in inner cores and
\emph{collisionless} in outskirts. The categorical reading is different.
There is one $D$-sector readout; it produces two regime-separated
nonlinearities at the cluster-substructure scale, organized by a single
length scale $\ellsub$ that depends on subhalo mass, not on a particle
parameter. Inside $\ellsub$ the readout is dominated by inner promotion of
the boundary-holonomy mass $\Mpa(R)$ in the host capacity well. Outside
$\ellsub$ the capacity-screening factor relaxes to unity and the
Newtonian-limit theorem governs.

\paragraph{One readout, not two.}
The categorical position is that the dichotomy is an \emph{artefact} of
asking the question in particle-DM language. In Category $\tau$ there is
no second component, no density-dependent cross-section, no hybrid; the
inner and outer phenomenologies are two regime-projections of one
underlying readout (V.D121 capacity skeleton + V.D272 boundary-holonomy
mass closure + V.D261 cluster screening). The three subsections below
make this concrete.

\subsection{Inner regime: \texorpdfstring{$\Mpa(R)$}{M\_partial(R)} inner-promotion}\label{subsec:inner}

\paragraph{Mechanism.}
At cluster radii $R/\Rtwohun \lesssim 0.2$, the host capacity gradient is
maximal and the boundary-holonomy mass $\Mpa(R)$ is enhanced relative to
its outskirts value (V.D272~\readout). The capacity-skeleton population of
sub-threshold peaks is thereby \emph{preferentially promoted} above the
spectroscopic luminosity threshold inside this projected radius. The
result is an excess in the projected member-galaxy count at
$R/\Rtwohun \lesssim 0.2$ relative to a $\Lambda$CDM expectation that
carries no $\Mpa$ at all (the TNG-Cluster suite,~\cite{NelsonTNGC2024},
solves only the standard $\Lambda$CDM Einstein--Boltzmann system; the
boundary-holonomy contribution is structurally absent).

\paragraph{Why this is sign-correct.}
$\Lambda$CDM substructure suffers numerical disruption inside $\sim 0.2
\Rtwohun$ even after Hopkins-style force-resolution-free reanalysis
(see~\cite{Natarajan2026} App.~B and~\cite{Hopkins2023,
ChiangTNGC2026}); the resulting simulated radial distribution carries
a residual inner \emph{dearth}, not an inner excess. The Category-$\tau$
direction is opposite: $\Mpa(R)$ inner-promotion concentrates the readout
inward. The amplitude is not derivable from existing primitives
($\Mpa(R)$ functional form deferred to a future Book V cluster-interior
module); the \emph{direction} is structurally guaranteed by V.D272
boundary-holonomy enhancement at the cluster nodal-capacity maximum.

\paragraph{V.T203 conservation.}
The V.T203 conjectural status~\cite{KSZCategorical} L811--814
(cluster capacity mass discrepancy, $D_\tau \sim 2$--$4$ vs.\
$D_{\mathrm{obs}} \sim 5$--$7$) governs the \emph{integrated}
cluster-scale mass closure, not the projected radial distribution of
peaks. The present~\S\ref{subsec:inner} therefore neither narrows nor
inherits V.T203, but \emph{relocates} its empirical signature: the
mass discrepancy can in principle be partially absorbed by the
boundary-holonomy enhancement near the cluster centre, which has not
been previously quantified.

\subsection{Outer regime: Newtonian-limit on tidally-pruned capacity wells}\label{subsec:outer}

\paragraph{Free consistency check.}
At radii $r \sim 10$--$100\,\kpc$, the cluster-scale screening factor is
\[
  1 + \frac{r}{\elltauv} \;=\; 1 + \frac{r}{c/(H_0 \sqrt{1 - \imark})}
  \;\approx\; 1 + (1.6\text{--}16)\times 10^{-5}
\]
(using the cosmological identity $\elltauv = c/(H_0\sqrt{1-\imark})
\approx 5.48\,\mathrm{Gpc}$~\cite{KSZCategorical}, V.D261).%
\footnote{The V.D261 entry in the corpus historically carried a unit-label
slip (Mpc vs.\ Gpc) documented in~\cite{KSZCategorical}~\S{}6; the present
note adopts the corrected Gpc-scale identity consistently. The dressed value
incorporating the V.T207 $\sqrt{1-\imark}$ pre-factor lies in the range
$5.48$--$13.0\,\mathrm{Gpc}$; the figure quoted here is the bare value
appropriate for the screening-factor derivation. The $K_0(r/\elltauv)$
near-field profile of V.T85 is the Yukawa-Stokes signature of a massive
propagator ($\eta_\tau = 0$, mean-field universality class), not a marginal-RG
inheritance --- a distinction sharpened in the Phase 13 Path C${}'$ synthesis
and disjoint from cousin-paper marginality results.}
The Roche-like tidal expression for subhalo truncation
(Book~V~ch.36~\readout) therefore reduces to the standard Newtonian form
to one part in $\sim 10^{5}$, and the predicted tidal-truncation-radius
distribution is collisionless-CDM-like \emph{by construction}.
Natarajan's Figure~3 (the SIDM $r_t$ distribution sits nearly an order of
magnitude below the CDM-like observed band; the specific median
$r_t^{\mathrm{SIDM}}/r_t^{\mathrm{CDM}} \approx 0.73$ at the $5\sigma$
rejection level is reported in Chiang+2026a~\cite{ChiangTNGC2026}) is
therefore matched without invoking dark-sector self-interactions; this
is a consistency check on the dichotomy's \emph{outer} side, not a
derivation of the specific $r_t$ value.

\paragraph{What this does and does not commit.}
This subsection commits the framework only to \emph{statistical
consistency with CDM at outer subhalo radii}, which is identical to the
collisionless-CDM expectation that Natarajan~et~al.\ find satisfied. It
does not commit a value of $r_t$ for any specific subhalo, nor an
amplitude of the $r_t$ distribution width. Those follow from the host
density profile and orbital-pericentre distribution, which are
external inputs to the present~\readout.

\subsection{Transition scale \texorpdfstring{$\ellsub$}{l\_cl,sub}}\label{subsec:transition}

\paragraph{Closed-form prediction.}
The transition between the inner $\Mpa$-dominated regime
(\S\ref{subsec:inner}) and the outer Newtonian-limit regime
(\S\ref{subsec:outer}) is governed by the
\emph{classical-validity length applied to the member subhalo's own mass}:
\begin{equation}\label{eq:ellsub}
  \ellsub(\Msub) \;=\; \sqrt{\frac{G \Msub}{\aMOND}},
  \qquad \aMOND \;=\; \frac{c H_0 \imark}{2}\;\text{(V.D232)},
\end{equation}
where $\imark$ is the master $\tau$-constant (Book~V V.D232 in
the capacity-gradient framework~\cite{KSZCategorical}). The host-cluster
$\elltauv$ (Gpc scale; see above footnote) and the host-cluster classical-validity scale
$\ellcl(\Mtwohun) \approx 1.2\,\Mpc \approx 0.45\,\Rtwohun$ are both
\emph{outside} the inner transition; the empirically relevant length
in the cluster-substructure context is the \emph{member} scale
$\ellsub(\Msub)$, not the host scale.

\paragraph{Numerical anchor.}
For $\Msub \in [10^{10}, 10^{11}]\,\Msun$ (the dominant mass range of
the Natarajan spectroscopically-confirmed member sample as tabulated
in~\cite{Natarajan2026}~\S\S{}2--3 and Fig.~1; the underlying SHMF
extends to $\sim 10^{13}\,\Msun$, but the spectroscopic and dynamical
weight in the GGSL-relevant cluster-member population sits at the lower
end),
\[
  \ellsub \;\in\; [3.4,\; 10.8]\,\kpc,
\]
in order-of-unity agreement with the empirical break at $r \sim 10\,\kpc$
across all three MACS clusters (MACS~J0416, J1206, J1149). The match is
\emph{categorical} (the existence and location of a transition scale,
not its precise value): Category $\tau$ predicts a transition, $\Lambda$CDM
and MOND predict none, and the observed break sits in the predicted
band. The host concentration ($c$, $r_s$) dependence of $\ellsub$ is
deferred to a closed-form derivation via the full V.D259 nonlinear
capacity equation (a future Book~V cluster-interior module); the
present order-of-unity match is robust to the host-$c$ variation across
$c \in [3.5, 8]$ at the level of $\pm 0.3\,$dex.

\paragraph{Distinguishability of the transition scale.}
The closed-form~\eqref{eq:ellsub} is structurally a MOND-proxy at the
\emph{subhalo} (member-galaxy) scale, where MOND traditionally works.
This is \emph{not} the host-cluster scale where MOND fails the
missing-mass test (Famaey, Pizzuti \& Saltas~\cite{Famaey2025},
Fig.~4 of which shows residual missing mass in CLASH-sample relaxed
clusters with central density ratio $\sim 10\times$ the gas density at
inner $\Mpc$). The Category-$\tau$ inheritance of MOND-proxy is therefore
regime-bounded:
\begin{itemize}[leftmargin=1.4em,nosep]
  \item At subhalo $\Msub \sim 10^{10}$--$10^{11}\,\Msun$:
        $\aMOND$-scale validity holds (V.P68 + V.R175~\cite{KSZCategorical}); $\tau$ inherits the MOND-proxy
        transition radius.
  \item At host-cluster $\Mtwohun \sim 10^{15}\,\Msun$:
        the validity regime is exited; the cluster-scale failure mode
        of MOND~\cite{Famaey2025} is recovered into V.T203's
        conjectural slot, not silently absorbed.
\end{itemize}
\noindent
SIDM core-collapse predicts $r_{\mathrm{trans}}^{\mathrm{SIDM}} \propto
(\sigma/m)^{\alpha}$ with no scaling against $\Msub$ at fixed cross-section;
Category $\tau$ predicts $\rtran \propto \Msub^{1/2}$. The
\emph{host-mass independence} of $\rtran$ at fixed subhalo mass function
is the headline pre-registered falsifier (\falsifierID{3} of
\S\ref{sec:falsifiers}).


% -----------------------------------------------------------------------
\section{Consistency checks passed}\label{sec:consistency}

Three orthogonal Natarajan diagnostics admit Category-$\tau$ consistency
readings at the \readout{} level.

\begin{description}[leftmargin=1.7em,itemsep=0.4em]
\item[\S{}4.1 Outer $r_t$ at $10$--$100\,\kpc$.]
The screening factor of \S\ref{subsec:outer} reduces the Book~V ch.36 Roche
expression to the standard Newtonian form at one part in $\sim 10^{5}$.
The Natarajan Figure~3 SIDM-vs-CDM distribution (with the Chiang+2026a
companion's $r_t^{\mathrm{SIDM}}/r_t^{\mathrm{CDM}} \approx 0.73$ median
at the $5\sigma$ rejection level~\cite{ChiangTNGC2026}) is matched by
construction. A consistency check, not a derivation of any specific $r_t$.

\item[\S{}4.2 Inner member-number excess at $R/\Rtwohun \lesssim 0.2$.]
The boundary-holonomy mass $\Mpa(R)$ inner-promotion mechanism
(\S\ref{subsec:inner}) produces a sign-correct concentration at small
projected radii. The Natarajan Figure~2 $5\sigma$--$40\sigma$ excess is
matched in direction; amplitude is deferred (G3 in~\S\ref{sec:gaps}).
The F1-wave $\Mpa(R)$ structural-closure analysis (Phase~12, V.D352 +
V.T348~\cite{Natarajan2026} cross-reference) predicts a $\tau$-side
inner-excess band consistent with the \emph{inner half} of the Natarajan
band ($\sim 5\sigma$--$20\sigma$), with the high-tail residual
($\gtrsim 30\sigma$) plausibly absorbing into the slope-axis
contribution that the V.D272 inner-promotion mechanism cannot supply.
TNG-Cluster~\cite{NelsonTNGC2024} cannot reproduce this structurally
because the underlying Einstein--Boltzmann system carries no
$\Mpa(R)$ contribution; the $5\sigma$--$40\sigma$ residual survives
the $10^6$-iteration Hopkins-style bootstrap correction in
Natarajan~Appendix~B~\cite{Natarajan2026,Hopkins2023}, so the
force-resolution-loss escape hatch is already closed by the source
analysis before the $\tau$-mechanism is invoked. Count-axis and
slope-axis contributions are partially degenerate at present
single-cluster precision; the count-vs-slope joint inference roadmap
is~\S\ref{sec:gaps} G1.

\item[\S{}4.3 Transition scale $\rtran \approx 3.4$--$10.8\,\kpc$.]
The closed-form $\ellsub$ of~\eqref{eq:ellsub} matches the empirical break
at $r \sim 10\,\kpc$ across the three MACS clusters to order unity, for
the $10^{10}$--$10^{11}\,\Msun$ subhalo mass range that dominates the
spectroscopically confirmed sample. The Natarajan~\S{}7 inner-density
transition and~\S{}8 outer truncation-radius transition share this
single length scale.
\end{description}


% -----------------------------------------------------------------------
\section{Pre-registered falsifiers (N16 surface)}\label{sec:falsifiers}

We register the following falsification surface (\falsifierID{1}--\falsifierID{11})
as the cluster-substructure $D$-sector dichotomy entry in the Panta Rhei
falsification ledger, continuing the N9--N15 sequence from
\cite{KSZCategorical} (N12--N14) and \cite{LRDCategoricalv33} (N15).
Rows \falsifierID{1}--\falsifierID{8} were registered in the original
research wave (R5 falsification-surface refinement); \falsifierID{9}
(joint $\imark$ universality) and \falsifierID{10} ($\Mpa(R)$ radial
shape) were added at Phase~9 (Workstream~C); and \falsifierID{11}
(concentration-dependent saturation radius $R_\star(c) = 1.25\,\Rtwohun/c^2$)
was added at Phase~12 from the V.D352 + V.T348 Chern--Weil partial-domain
holonomy structural closure.
Each row is \evstag{} and is testable with named near-term surveys
(Euclid, Rubin~LSST, JWST cluster-lensing follow-up, FRBs, SKA, Roman,
CMB-S4). The headline row is \falsifierID{3}; the consistency floor is
\falsifierID{1}.

\begin{table*}[!t]
\footnotesize
\centering
\renewcommand{\arraystretch}{1.15}
\caption{N16 falsification surface — cluster-substructure $D$-sector
dichotomy entry. Rows \falsifierID{1}--\falsifierID{8} originate from the
R5 falsification-surface refinement; \falsifierID{9}--\falsifierID{10}
were added at Phase~9 (Workstream~C); \falsifierID{11} was added at
Phase~12 from the V.D352~+~V.T348 Chern--Weil partial-domain holonomy
structural closure. Every row is \evstag{}.}
\label{tab:n16}
\begin{tabular}{@{}p{0.6cm} p{4.8cm} p{3.6cm} p{1.6cm} p{2.6cm} p{1.7cm}@{}}
\toprule
ID & Prediction & $\tau$ & CDM & SIDM core-collapse & MOND \\
\midrule
\falsifierID{1} & Outer $r_t$ at 10--100~$\kpc$ NFW-Roche-like
  & match\,(\S{}4.1)
  & match
  & $\sim 10\times$ smaller
  & ambiguous \\
\falsifierID{2} & Inner $R/\Rtwohun \lesssim 0.2$ \emph{number} excess
  & match\,(\S{}4.2); $5$--$7\times$, $5\sigma$--$20\sigma$ band
  & dearth (TNG-C)
  & dearth
  & no prediction \\
\falsifierID{3} & \textbf{Host-mass independence of $\rtran$ at fixed SHMF}
  & $\rtran \propto \Msub^{1/2}$ \emph{(headline)}
  & no $\rtran$
  & $\propto (\sigma/m)^{\alpha}$
  & wrong scale \\
\falsifierID{4} & Transition $\rtran \approx \ellsub \approx 3.4$--$10.8~\kpc$
  & match\,(\S{}4.3, \eqref{eq:ellsub})
  & no $\rtran$
  & $\Msub$-independent
  & host-scale, not subhalo-scale \\
\falsifierID{5} & SHMF slope $dn/d\Msub \propto \Msub^{-1.9}$
  & slope $\sim -2$ (capacity peaks);
    $A_\tau/A_{\rm TNG-C} \in [0.5, 1.5]$ per V.T213 floor
  & match
  & suppressed at low $\Msub$
  & no prediction \\
\falsifierID{6} & Cluster-scale missing-mass profile within central 1~Mpc
  & V.T203 conjectural slot~\cite{KSZCategorical}
  & full match
  & residual
  & $\sim 10\times$ gas, inner core~\cite{Famaey2025} \\
\falsifierID{7} & FRB scintillation: cluster sightlines through $\Mpa$ peaks
  & dispersion-measure bumps where $\Mpa(R)$ peaks
  & smooth
  & smooth
  & no prediction \\
\falsifierID{8} & GGSL probability scaling with $\Msub$
  & member-density dominated
  & low (dearth)
  & slope-dominated (cusp)
  & no prediction \\
\falsifierID{9} & \textbf{Joint $\imark$ universality across LRD + cluster samples}
  & shared pre-factor $\imark = 2/(\pi+e) \approx 0.341$ at both
    LRD aperture-saturation and cluster $\Mpa(R)$
  & no relation
  & no relation
  & no relation \\
\falsifierID{10} & Radial-shape discriminator on $\Mpa(R)$
  & shape between $1/R$ and $1/R^3$ inside $\Rtwohun$
  & no $\Mpa$
  & no $\Mpa$
  & no $\Mpa$ \\
\falsifierID{11} & Concentration-dependent saturation radius
    $R_\star(c) = 1.25\,\Rtwohun/c^2$
  & $R_\star/\Rtwohun \in [0.02, 0.10]$ for $c \in [3.5, 8]$;
    derived from NFW partial-domain Chern--Weil
  & no $\Mpa$
  & no $\Mpa$
  & no $\Mpa$ \\
\bottomrule
\end{tabular}
\end{table*}

\paragraph{Headline: \falsifierID{3}, host-mass independence of $\rtran$.}
The discriminating prediction is the cluster-mass dependence of the
transition radius. Category $\tau$ predicts $\rtran \propto \Msub^{1/2}$
with no explicit host-mass scaling at fixed subhalo mass function; SIDM
predicts $\rtran \propto (\sigma/m)^\alpha$ (subhalo-mass-independent at
fixed cross-section); CDM predicts no transition at all; MOND predicts
the wrong scale at the host mass.
The test requires a sample of $\gtrsim 30$ massive lensing clusters
with strong+weak lensing and substructure-resolved inner profiles
(consistent with the Meneghetti~+2020 / Despali~+2017 survey-forecast
standard~\cite{Meneghetti2020}; a tighter band-comparison $\chi^2$
formulation may reduce this to $\sim 20$ in a future power-calculation
pass, but the conservative figure stands here); \falsifierID{9} adds
$\sim 50$ SLACS-class galaxy-scale lenses plus $\sim 30$ LRD lenses on
top, with the shared $\imark$ parameter delivering a tighter joint
constraint than any single survey delivers alone. The Euclid wide
survey plus Rubin LSST Year-2~\cite{EuclidRedbook2024,RubinLSST2025}
deliver the first sub-sample by 2027--2029; the full headline test
also requires JWST + Roman kpc-resolved $\gamma(r)$ stacking on the
2028--2030 horizon, not Euclid+LSST alone.

\paragraph{Consistency floor: \falsifierID{1}, outer $r_t$.}
The Roche-like tidal extents at $r \sim 10$--$100\,\kpc$ are matched by
construction in Category $\tau$ (\S\ref{subsec:outer}); the Natarajan
inference here is therefore \emph{consistent}, not predictive. This row
survives even if all other R-wave outcomes degrade to negative results.

\paragraph{V.D272 / V.T213 coupling disclosure and $\imark$ joint-prediction lock.}
\falsifierID{3} and \falsifierID{6} share their boundary-holonomy mass
anchor (V.D272) with~\cite{LRDCategoricalv33}~Signature~3, which
attaches the same anchor to aperture-dependent saturation at LRD scales.
The coupling is \emph{quantitative}: both the LRD signature-3 $6.65$
saturation ceiling and the cluster-inner $\Mpa(R)$ amplitude in
\falsifierID{2}/\falsifierID{4} depend on the same $\tau$-master-constant
pre-factor $\imark = 2/(\pi+e) \approx 0.341$ via V.D272. The two predictions
are therefore \emph{not independently falsifiable}: a single $\imark$
parameter generates two empirical signatures on different observables,
giving a positive consistency check at the cost of joint exposure
(\falsifierID{9}). Any null result on the LRD-side SLACS-saturation
prediction propagates immediately to a tension on the present
$\Mpa(R)$-based inner-promotion; conversely, any aperture-independent
$\Mpa(R)$ measurement would falsify both predictions jointly and force
a structural revision of V.D272. This coupling is disclosed here per
the R4 conservatism ledger and the F5 falsification-surface refinement.

\paragraph{Row-independence audit.}
The eleven N16 rows span approximately five independent physical axes,
not eleven. The intra-N16 couplings are: \falsifierID{3} $\leftrightarrow$
\falsifierID{4} (both anchored on $\rtran = \ellsub$; a null on
\falsifierID{4} cascades to \falsifierID{3}); \falsifierID{2}
$\leftrightarrow$ \falsifierID{7} $\leftrightarrow$ \falsifierID{10}
(all anchored on $\Mpa(R)$ inner-promotion; a null on \falsifierID{2}
propagates to \falsifierID{7} and \falsifierID{10}); \falsifierID{11}
$\leftrightarrow$ \falsifierID{10} (both probe $\Mpa(R)$ radial shape).
The shared $\imark$ pre-factor across \falsifierID{2}, \falsifierID{4},
\falsifierID{6}, \falsifierID{9} (V.D272 + LRD-S3 joint lock above) is
the cross-program coupling. Net: the eleven rows resolve to roughly
\emph{five independent axes} — transition scale, $\Mpa(R)$ inner-promotion,
SHMF slope/amplitude, V.T203 cluster mass discrepancy, and FRB
scintillation. The shared $\imark$ joint exposure is, equivalently, a
\emph{joint parameter constraint}: a future $\sim 50$-SLACS +
$\sim 30$-LRD + $\sim 30$-cluster joint fit on $\imark$ is a tighter
single-parameter constraint than any single survey delivers.

\paragraph{Timeline tiers.}
First decisive data per row: 2026--2028 (JWST cluster-lensing follow-up
of Frontier-Fields and CLASH targets via the BUFFALO, MAGNIF, CANUCS,
and GLASS-JWST programs); 2027--2029 (Euclid wide survey plus Rubin
LSST Year-2 first sub-sample); 2028--2030 (full Euclid + LSST + JWST +
Roman + FRB cluster sightlines); 2035+ (SKA, Roman wide,
CMB-S4).


% -----------------------------------------------------------------------
\section{Honest gaps and the structurally honest tension}\label{sec:gaps}

The wave that produced this note (\texttt{research-wave/02-R1\ldots R5})
explicitly closed two of the three binding gaps from the
prior red-team triage at the partial level (direction-and-mechanism),
and produced one numerical anchor (\eqref{eq:ellsub}). It also returned
one principled negative finding that we now reframe as a falsifier.

\paragraph{G1 — inner $\gamma$ value (R2 NEGATIVE; reframed).}
The natural candidate mechanism for steep inner cusps in
Category $\tau$ was the \emph{saturation family} of Book V
(V.T36, V.D57, V.P16, V.T117). On re-inspection the family is
\emph{structurally anti-cusp}: V.P16 explicitly replaces every
configuration that would otherwise produce a general-relativistic
density singularity with bounded density. The Schwarzschild-radius
scale $\Rs(\Msub \sim 10^{11}\,\Msun) \sim 10^{-8}\,\kpc$ is 4--7
decades below Natarajan's target $r \lesssim 0.01\,\Rtwohun \approx
1$--$10\,\kpc$, so the saturation-replaces-singularity construction is
geometrically inert at the relevant scale. The V.T117 saturation
radius is a \emph{cosmological motif-saturation} length
($\sim 100\,\Mpc$ scale), not a kpc-scale halo-core radius.
The V.D305 fiber-holonomy primitive cited by triage agent A1 as a
candidate cross-context primitive turns out to be the BBN $^7$Be
radiative-capture correction (module
\texttt{BookV.Cosmology.BBNNuclearNetwork}), not a saturation
primitive; the actual Wave-R7 LRD-1 heavy-seed derivation uses
V.D-LRD-1d, V.T88, V.T108, V.T110 plus $\Pi_{\mathrm{coh}}$
(\cite{LRDCategoricalv33}), all rotational/threshold mechanisms with
no density-slope crossover.
\smallskip\noindent
The categorical conclusion is that
\emph{Category $\tau$ predicts \textbf{flatter}, not steeper, inner
density slopes at $r \lesssim \kpc$ in cluster subhalos}, on the
\emph{opposite} side of Natarajan's parametric $\gamma \gtrsim 2.5$
dPIE inference (whose direct empirical anchor is Dutra et~al.\
2025~\cite{Dutra2025}, $\gamma \sim 2.9$ across three relaxed MACS clusters)
from SIDM core-collapse.
\smallskip\noindent
\textbf{Reframe as a conditional discriminating observation
(\emph{not} a dissolution of G1).} The GGSL probability $\Pggsl$ is
sensitive to the inner mass distribution \emph{integrated} along the
line of sight; it carries contributions from \emph{two distinct axes}
that are partially separable in joint GGSL+kinematics inference (the
prototype technique is Granata et~al.\ 2022~\cite{Granata2022}'s MUSE+HST
joint strong-lensing + member-velocity-dispersion modelling on
MACS~J0416, with current state-of-the-art HFF/CLASH lens models in
Bergamini et~al.\ 2023~\cite{Bergamini2023}):

\begin{description}[leftmargin=1.7em,itemsep=0.3em]
\item[Count axis.] The number and spatial concentration of
member-galaxy subhalos providing GGSL targets. Category $\tau$
addresses this axis sign-correctly via the $\Mpa(R)$ inner-promotion
mechanism (\S\ref{subsec:inner}).
\item[Slope axis.] The per-subhalo inner density profile $\gamma$. SIDM
core-collapse addresses this axis by steepening; Category $\tau$
predicts \emph{flatter}, not steeper, slopes via V.P16. Importantly,
Tokayer~et~al.~(2024)~\cite{Tokayer2024} already exhausted the
\emph{count-axis-only} approach via intra-cluster baryonic mass
rearrangement and found it insufficient to match the observed
$\times 3.5$--$\times 8$ $\Pggsl$ boost in Natarajan~\S{}7. Thus the
$\tau$-readout cannot dissolve G1 as a category error; the slope axis
is empirically required to be \emph{somewhere}, and the conditional
question is where.
\end{description}

\noindent
The discriminating test, as a conditional:
\begin{quote}\itshape\readout
At fixed inferred $\Pggsl$, does the GGSL+inner-mass-content joint
inference attribute the signal predominantly to count axis ($\tau$ via
$\Mpa(R)$ inner-promotion) or to slope axis (SIDM core-collapse with
$\gamma \gtrsim 2.5$)? If predominantly count axis,
Category $\tau$ matches without invoking a steep inner slope; if
predominantly slope axis, $\tau$ requires the deferred metric-sector
progenitor of \S\ref{sec:gaps} (V.P144 + V.D123 trace-projection;
program-level open item) before answering. Galaxy-scale corroboration of
core-collapsed compact substructure is independently emerging outside
the cluster regime (Vegetti et~al.\ 2026 on JVAS~B1938+666 via
VLBI~\cite{Vegetti2026}; Powell et~al.\ 2025 via ALMA~\cite{Powell2025}),
which Natarajan~et~al.\ \S{}10 cite as cross-context evidence; the
present note remains cluster-scale-only. Future JWST cluster-lensing
+ kinematic IFU follow-up will discriminate (BUFFALO, MAGNIF, CANUCS,
GLASS-JWST programs are the near-term anchor; $\sim 2028$--$2030$ horizon).
\end{quote}

\paragraph{G2 — SHMF amplitude (R1 PARTIAL).}
The slope $\sim -2$ of the cluster member-galaxy capacity-peak abundance
follows from peak statistics of the CMB/BAO-constrained Gaussian random
field over the underlying capacity field (R1 derivation, ~\readout); the
recovered slope agrees with the Natarajan-observed $dn/d\Msub \propto
\Msub^{-1.9}$ to within slope uncertainty. The \emph{amplitude} of the
SHMF is not derivable from existing primitives. Flagged for future
Book V cluster-interior scaffolding.

\paragraph{G3 — inner-excess amplitude (R1 PARTIAL).}
The sign-correct inner concentration of member counts at $R/\Rtwohun
\lesssim 0.2$ via $\Mpa(R)$ inner-promotion (\S\ref{subsec:inner})
is structurally guaranteed; the $5\sigma$--$40\sigma$ \emph{amplitude}
is not derivable from existing primitives. Flagged for future Book V
scaffolding.

\paragraph{Pre-registration discipline.}
The three gaps above are \emph{not} promised against existing primitives.
A future Book V module --- tentatively ``Cluster Member Distribution from
the Capacity Skeleton with Boundary-Holonomy Inner Promotion'' --- would
close G2 and G3 simultaneously; closing G1 requires a genuinely new
\emph{boundary-holonomy radial-steepening} mechanism distinct from both
V.D272 inner-promotion (which is a number-density mechanism) and the
saturation family (which is structurally anti-cusp). Such mechanism, if
it exists, would lift this note to the deferred F3 form factor. As of
the present authoring it does not exist in the published $\tau$-corpus,
and the structurally honest position is to commit to G1's reframe-as-falsifier rather than to scaffold a new primitive on demand.


% -----------------------------------------------------------------------
\section{Distinguishability statement}\label{sec:distinguish}

Category $\tau$ is structurally distinct from every active alternative on
the Natarajan landscape:

\begin{description}[leftmargin=1.7em,itemsep=0.3em]
\item[CDM (collisionless).] Same outer $r_t$ prediction by construction
(\S\ref{subsec:outer}); different inner mechanism --- CDM has no
$\Mpa(R)$ and predicts the observed \emph{dearth}, not excess; no
transition scale.
\item[Constant-$\sigma$ SIDM.] Different mechanism --- saturation
(anti-cusp) vs.\ collisional steepening; opposite $\gamma$-direction
(\S\ref{sec:gaps} G1); transition scale parameter dependence:
$(\sigma/m)$ vs.\ $\Msub^{1/2}$ (\falsifierID{3}).
\item[Velocity-dependent SIDM.] Same as constant-$\sigma$ on the
$\Msub$-vs-$\sigma/m$ scaling discriminator \emph{up to a residual
$v_{\rm disp}(\Msub)$ running} via the Rutherford-like cross-section of
Yang \& Yu 2022~\cite{YangYu2022} adopted in Nadler et~al.\ 2025; this
running is much weaker than $\Msub^{1/2}$ and carries the
\emph{opposite sign} in the mildly-collisional regime (smaller subhalos
have lower $v_{\rm disp}$, hence larger effective $\sigma/m$, hence
\emph{earlier} core-collapse and smaller $\rtran$). The mildly-collisional
long-lasting core-phase mechanism of Errani et~al.\
2023~\cite{Errani2023} is what closes Natarajan~\S{}8's vel-dep escape
hatch. The \falsifierID{3} discriminator therefore survives against
vel-dep SIDM, but the foil is not identical to constant-$\sigma$.
Same as constant-$\sigma$ on the cusp/anti-cusp $\gamma$-direction
discriminator. Outer $r_t$-truncation tighter than CDM in vel-dep
cross-section models~\cite{Nadler2025}; Category $\tau$ predicts CDM-like
by construction.
\item[MOND.] V.T78 force-law exponent $n = 2$ vs.\ MOND $n = 1$;
$\aMOND$-scale validity regime-bounded
(\S\ref{subsec:transition}): inherited at subhalo $\Msub \sim
10^{10}$--$10^{11}\,\Msun$, exited at host $\Mtwohun \sim 10^{15}\,\Msun$
where MOND fails the cluster-mass test~\cite{Famaey2025}. Category $\tau$
does \emph{not} invoke a universal $\aMOND$ across all scales; the
$\ellsub$ length \emph{is} a $\sqrt{G\Msub/\aMOND}$ expression but with
$\aMOND$ a derived constant from the cosmological capacity-gradient
framework (V.D232, V.P68, V.R175~\readout), not a postulate.
\item[Hybrid CDM+SIDM.] Ontologically distinct --- one $D$-sector, no
second component, no density-dependent cross-section
(\S\ref{sec:two-regimes}). Observationally distinct on
\emph{host-concentration} dependence of $\rtran$: hybrid CDM+SIDM
still carries a density-dependent transition (slope axis activates
inside $\rho_{\rm crit}$), giving a $\rtran$ that depends on host
concentration $c$ at fixed $\Msub$; Category $\tau$'s $\rtran \propto
\Msub^{1/2}$ scaling carries no explicit host-$c$ dependence. The two
are discriminable in a $\sim 30$-cluster sample across the
$c \in [3.5, 8]$ range of relaxed CLASH/HFF cluster lenses.
\item[Verlinde / MOG.] Different mechanism class --- Category $\tau$
inherits the Newtonian $n = 2$ force law from V.T78~\cite{KSZCategorical},
not entropic gravity or Yukawa-modified gravity; no parametric
free parameter to tune against cluster lensing.
\end{description}

\noindent
The minimum testable distinguishability set within the
\falsifierID{1}--\falsifierID{11} surface (\S\ref{sec:falsifiers}) is
\falsifierID{3} (Category $\tau$ uniquely predicts host-mass
independence of $\rtran$ at fixed SHMF) plus \falsifierID{2}
($\Mpa(R)$ inner-promotion produces excess; all four alternatives
produce dearth or no prediction). Combined, these two rows form a
minimum-cluster-sample-size discriminator of $\sim 30$ massive
clusters at strong+weak lensing precision, which Euclid+LSST will
deliver by $\sim 2029$.


% -----------------------------------------------------------------------
\section{Status discipline}\label{sec:status}

Every $\tau$-claim in this note is tagged \readout{} or \evstag{}; no
\textsc{[formal]} or \textsc{[hinge]} claim is made on the Natarajan
empirical observables. The section-to-tag mapping is given in
Table~\ref{tab:status-mapping}.

\begin{table}[!h]
\footnotesize
\centering
\renewcommand{\arraystretch}{1.15}
\caption{Status-tag mapping (\readout{} = consistency reading / interpretation;
\evstag{} = external-validation seam / falsifier).}
\label{tab:status-mapping}
\begin{tabular}{@{}p{2.6cm} p{3.3cm} p{0.9cm}@{}}
\toprule
Section & Claims & Tag \\
\midrule
\S\ref{subsec:inner} ($\Mpa(R)$ inner-promotion)
  & sign-and-mechanism reading & \readout \\
\S\ref{subsec:outer} (outer $r_t$ via Newtonian limit)
  & consistency check & \readout \\
\S\ref{subsec:transition} (transition $\ellsub$)
  & numerical reading & \readout \\
\S\ref{sec:consistency} (three consistency checks)
  & all \S{}4.x & \readout \\
\S\ref{sec:falsifiers} (N16.1--N16.11)
  & future-observation falsifiers & \evstag \\
\S\ref{sec:gaps} (G1 reframed; G2, G3 flagged)
  & pre-registered gaps + reframed falsifier & \evstag \\
\S\ref{sec:distinguish} (foil paragraphs)
  & positioning & \readout \\
\bottomrule
\end{tabular}
\end{table}

\noindent
V.T203 (Cluster Capacity Mass Discrepancy) inherits its
\texttt{skeleton}/conjectural status from~\cite{KSZCategorical}
L811--814 verbatim; this note neither narrows nor extends V.T203, only
structurally relocates its empirical signature
(\S\ref{subsec:inner} \emph{V.T203 conservation} paragraph). The V.T213
five-cluster universality is inherited as a calibration prerequisite on
integrated cluster mass, not as a substructure-resolved theorem
(\S\ref{sec:tau-position} \emph{KSZ cluster-scale precedent} paragraph).


% -----------------------------------------------------------------------
\section{What this note does NOT claim}\label{sec:non-claims}

\begin{enumerate}[leftmargin=1.4em,itemsep=0.2em]
\item It does \emph{not} claim Category $\tau$ \emph{derives} the
$\gamma \gtrsim 2.5$ inner density slope, the SHMF amplitude, the
$5\sigma$--$40\sigma$ inner-excess amplitude, or the specific $r_t$
distribution width inferred by Natarajan~et~al.\ from any existing
$\tau$-corpus primitive. Each is flagged as a pre-registered gap
(\S\ref{sec:gaps}).
\item It does \emph{not} claim that $\tau$-physics is Lean-verified
as external physics. The cluster-substructure context is outside the
scope of the Lean-formalized $\tau$-internal-mathematics envelope
(see the dossier~\cite{BareMetalSpineV2}~App.~C).
\item It does \emph{not} claim an externally-settled resolution of the
small-scale CDM tensions or of any Clay Millennium problem related to
gravitational lensing or cluster physics.
\item It does \emph{not} predict the IllustrisTNG-Cluster
simulation~\cite{NelsonTNGC2024} as a $\tau$-readout; the simulations
are external $\Lambda$CDM-expectation comparators, not $\tau$-derivations.
\item This note is a privately-deposited arXiv response in the Panta Rhei
research-note format and is not a substitute for direct dialogue, which
the authors actively welcome and to which the present pre-registration
is offered as an opening contribution.
\item It does \emph{not} supersede the conjectural status of V.T203
in~\cite{KSZCategorical}; the V.T203 cluster mass discrepancy remains
an open conjectural item, structurally relocated but not narrowed
by the present note (\S\ref{subsec:inner}).
\end{enumerate}


% -----------------------------------------------------------------------
\section*{Acknowledgements}

We gratefully acknowledge P.~Natarajan, B.~T.~Chiang, and I.~Dutra for the
open-access deposit of~\cite{Natarajan2026}, for the rigour of the
$10^6$-iteration force-resolution bootstrap injection (their Appendix~B,
building on the Hopkins+2023 FIRE-3 benchmark~\cite{Hopkins2023}), and
for the structural foreclosure of the conventional SIDM resolution in
their \S{}10 (with the Chiang+2026b companion
analysis~\cite{ChiangSimResolution2026}); each of these methodological
commitments materially shaped the categorical reading developed here.
Engagement, criticism, and falsification of any of the eleven
pre-registered N16 rows would be welcomed and would advance the program.

% -----------------------------------------------------------------------
\bibliographystyle{unsrt}
\bibliography{references}

\end{document}
